\documentclass[a4paper,12pt]{extarticle}
\usepackage[utf8]{inputenc}
\usepackage[T2A]{fontenc}
\usepackage[russian]{babel}
\usepackage{geometry}
\geometry{left=20mm, right=15mm, top=20mm, bottom=20mm}
\usepackage{tempora}
\usepackage{enumitem}
\usepackage{parskip}
\usepackage{titlesec}
\titleformat{\section}{\Large\bfseries}{}{0pt}{}
\titleformat{\subsection}{\large\bfseries}{}{0pt}{}
\setlength{\parindent}{0pt}

\begin{document}

\begin{center}
{\Large\bfseries Речь для защиты --- HS:BG RL}
\end{center}

\textbf{Тайминг}: $\sim$7 минут 10 секунд на 13 слайдов. Лимит КР --- 10 минут, рекомендация --- 7.

\textbf{Правила репетиции:}
\begin{enumerate}[noitemsep]
  \item Прочитай вслух с секундомером, засеки каждый слайд отдельно
  \item Если слайд выходит дольше --- сокращай. Не подгоняй сам себя в конце
  \item Ключевые слайды \textbf{не сокращай}: 6 (архитектура), 7 (главный результат)
  \item Запинки естественные, не звучи как диктор
  \item Обязательно тыкай пальцем в слайды 9--10 (attention) --- даёт комиссии якорь
\end{enumerate}

\hrulefill

\section*{Слайд 1 --- Титул (10 сек)}

Здравствуйте. Меня зовут Пономарев Тимур, группа БПМИ243. Тема --- применение обучения с подкреплением в играх жанра автобаттлеры.

\section*{Слайд 2 --- Мотивация (45 сек)}

Среда, на которой я работаю --- Hearthstone Battlegrounds. Это автобаттлер от Blizzard, более 20 миллионов игроков в месяц. Игра состоит из двух чередующихся фаз: \textbf{recruit} --- экономическая фаза, где игрок покупает юнитов, продаёт их, реролит магазин, апгрейдит таверну; и \textbf{combat} --- автоматический бой двух досок.

Почему эта среда интересна именно как RL-задача? Это богатый testbed: здесь одновременно есть POMDP, потому что доску оппонента ты видишь только в бою, стохастика в самом бою, комбинаторный action space, и ещё поверх всего экономическое планирование. Структурно эта среда даже богаче, чем Pokemon из недавней работы Guo на RLC 2025 --- там нет экономики и позиционирования. И важный момент --- жанр живой, Blizzard добавляет новые карты каждый сезон, поэтому бенчмарк растёт вместе со средой.

Цель работы --- разработать исследовательскую платформу и обучить агента.

\section*{Слайд 3 --- Формализация POMDP (35 сек)}

Формально это POMDP с gamma 0.99. Пространство наблюдений --- вектор размерности 1009, структурированный по зонам: глобальный контекст, доска, рука, магазин, discovery и информация об оппоненте. Каждая сущность кодируется 37 признаками. Пространство действий --- дискретное, 32 действия с динамическим маскированием нелегальных. Награда --- многокомпонентная с терминальным плюс-минус пять.

\section*{Слайд 4 --- Почему нетривиально + gap (45 сек)}

Почему задача нетривиальна: POMDP, стохастика, комбинаторика --- порядка десяти в восьмой конфигураций доски при 18 юнитах на 7 слотов, sparse reward. Credit assignment классический.

Что касается существующих сред: Fireplace поддерживает классический Hearthstone, но не Battlegrounds. Unity ML-Agents медленные из-за графики. HSReplay --- только реплеи без интерактивного управления. Ближайшие академические работы --- Leiden 2022 и LUT 2024 --- используют сильные упрощения: без аур, без Discovery, без позиционных эффектов. Поэтому для исследования я написал собственный движок с полной боевой системой.

\section*{Слайд 5 --- Симулятор (50 сек)}

Движок на Python, на момент отчёта --- около 3500 строк кода самого движка плюс 839 строк Gymnasium-обёртки. 11 модулей в пакете engine.

Ключевые подсистемы:
\begin{itemize}[noitemsep]
  \item \textbf{Entity System} с многоуровневым стеком баффов --- базовый, перманентный, ходовой, боевой, аурный слои
  \item \textbf{Event System} с приоритизированной очередью событий и детерминированной сортировкой триггеров
  \item \textbf{Combat} --- реализованы 10 ключевых механик, включая Deathrattle с цепными призывами, Reborn, Divine Shield, Cleave, Poisonous, Windfury, Magnetize, Avenge
  \item \textbf{Auras} с stateless пересчётом
  \item \textbf{Tavern} с триплетами, Discovery, заклинаниями
\end{itemize}

По охвату механик это наиболее полная реализация по сравнению с упомянутыми академическими работами.

\section*{Слайд 6 --- Архитектуры MLP vs Transformer (80 сек) --- КЛЮЧЕВОЙ}

Дальше --- архитектуры. Я сравнивал наивный MLP и Transformer-экстрактор.

Наивный MLP на плоском векторе $\mathbb{R}^{1009}$ плохо ложится на задачу по трём причинам: не различает семантические границы зон, не ловит взаимодействия между юнитами --- расовые синергии или связку Baron с Deathrattle-юнитами, и тратит параметры на паддинг-слоты переменной доски.

\textbf{Теперь как именно $\mathbb{R}^{1009}$ превращается в последовательность токенов для Transformer.} Наблюдение структурировано по зонам: 7 слотов доски, 10 слотов руки, 7 слотов магазина и 3 слота discovery. Это \textbf{27 entity-слотов}, каждый размерности 37. Каждый слот проходит через DecomposedEncoder и становится одним entity-токеном размерности d\_model, то есть 128. Дополнительно глобальный контекст из 7 признаков и информация об оппоненте из 3 конкатенируются и проходят через отдельный MLP, превращаясь в \textbf{обучаемый токен GLOBAL\_CTX} --- это 28-й токен. На выходе --- \textbf{последовательность $28 \times 128$}, которая идёт в стек из четырёх GTrXL-блоков и потом в PMA.

Теперь про компоненты, все из state-of-the-art работ:
\begin{itemize}[noitemsep]
  \item \textbf{DecomposedEncoder} из MARL-GPT кодирует разные группы атрибутов --- continuous, binary, race --- раздельными проекциями с суммированием
  \item \textbf{FiLM} из работы Perez 2018 --- мультипликативная модуляция токенов глобальным контекстом
  \item Токен \textbf{GLOBAL\_CTX} --- аналог CLS из BERT и scalar features из AlphaStar, позволяет картам обращаться к макро-контексту внутри attention
  \item \textbf{GTrXL-шлюзы} из работы Parisotto 2020 со смещением минус три --- ключ к стабильности: на старте шлюз почти закрыт, Transformer практически пропускает вход без изменений, агент учится простым вещам через линейные слои, а потом градиенты сами открывают шлюз
  \item \textbf{Multi-Seed PMA} из Set Transformer --- четыре обучаемых seed-вектора для агрегации множества карт переменного размера
\end{itemize}

Компоненты все известны из литературы, но сборка именно под автобаттлеры --- впервые.

\section*{Слайд 7 --- Sample Efficiency (85 сек) --- ГЛАВНЫЙ РЕЗУЛЬТАТ}

Главный экспериментальный результат. Слева --- реальные WandB-логи метрики board\_power, среднее значение силы доски. Оранжевая линия --- это Transformer, обучение остановлено на 300 тысячах шагов. Розовая --- MLP, 1.5 миллиона шагов.

Видно: Transformer выходит на плато board\_power около 15 примерно за 80 тысяч шагов, MLP достигает того же уровня за 400 тысяч. \textbf{То есть Transformer в пять раз быстрее по sample efficiency.}

Агент проходит три фазы обучения: сначала учится тратить золото, потом собирать синергии и триплеты, потом финальная шлифовка. Transformer проходит все три фазы в пять раз меньше шагов.

Про FPS --- тут надо честно сказать. Transformer считает 250 шагов в секунду, MLP --- 1200. Это per-step compute. Но wall-clock до плато у них уже примерно равен: у MLP 333 секунды, у Transformer 320. Почему так: MLP на плоской сети [512, 512] compute-bound на мелких матрицах, GPU у него простаивает, и добавление железа ему ничего не даст. Transformer же честно грузит GPU и масштабируется по batch size, по количеству устройств и по параметрам --- тут работают стандартные scaling laws. \textbf{С ростом compute Transformer будет становиться только быстрее MLP, а не просто равен.}

\section*{Слайд 8 --- PvP (35 сек)}

Второй эксперимент --- прямое столкновение Transformer против MLP. Три партии --- все три ушли в таймаут более 500 ходов. HP колебалось в районе 20--30. Ни один не смог добить другого.

Моя интерпретация: обе архитектуры \textbf{решили текущую конфигурацию среды}. При 18 картах и тирах 1--3 пространство стратегий слишком узкое, чтобы архитектура дала преимущество. То есть движок работает корректно, нет эксплойтов, но нужна более сложная среда --- и это мотивация для апдейтов после отчёта.

\section*{Слайд 9 --- Attention по слоям (40 сек)}

\textit{(тыкай пальцем/курсором в конкретные участки heatmap)}

Третий результат --- валидация через анализ attention-весов. Вот усреднённые карты внимания по слоям.

На \textbf{Layer 0} видно: внимание концентрируется \textbf{здесь} --- внутри доски, B0 по B6 --- и \textbf{вот здесь}, на токене GLOBAL\_CTX. Модель смотрит сама на себя и сверяется с контекстом.

А на \textbf{Layer 3} --- вот \textbf{здесь} появляются устойчивые cross-zone паттерны. Доска обращается к магазину. То есть модель на более глубоких слоях буквально задаёт вопрос <<какая карта из shop подошла бы моей доске>>.

\section*{Слайд 10 --- Attention heads (30 сек)}

\textit{(тыкай в каждую голову)}

Ещё интереснее --- если разложить Layer 0 по отдельным головам. \textbf{Head 0} концентрируется внутри доски. \textbf{Head 1} смотрит на магазин. \textbf{Head 3} --- глобальный контекст плюс Discovery. То есть головы \textbf{сами} специализировались на разных аспектах игрового состояния. Никто их к этому не принуждал --- PPO-градиент сам нашёл разделение труда.

\section*{Слайд 11 --- Reward + честность про ablation (40 сек)}

Отдельно про reward shaping. Это признанная открытая задача RL, см. Ng 1999 про potential-based shaping и Andrychowicz HER. Эмпирически я пробовал: сначала 5-компонентный shaping из отчёта, потом MC Oracle PBRS через C++ симуляцию исходов боёв, потом откат к sparse reward плюс tier-based curriculum. Значимого стабильного выигрыша ни один не дал. Это и есть мой частичный ablation по reward.

Про покомпонентный ablation архитектуры. \textbf{В рамках вычислительного бюджета данного исследования полная покомпонентная абляция нецелесообразна} --- фокус сделан на proof-of-concept архитектуры и валидации sample efficiency. Выбор компонент обоснован литературным обзором 60+ работ, документирован в репозитории.

\section*{Слайд 12 --- Апдейты + Future Work (30 сек)}

Коротко --- что сделано после сдачи отчёта. Пул карт расширен с 18 до 175, тиры 1--7. Combat engine портирован на C++ через pybind11 --- \textbf{270 тысяч боёв в секунду} на простых кейсах, 90 тысяч на сложных. Это в 40+ раз быстрее предыдущей C++-версии и примерно в \textbf{800 раз быстрее} чистого Python. Tier-based curriculum. Categorical Critic из DreamerV3. Кодовая база выросла с 3.5k до $\sim$10.7k Python плюс 3.7k C++ плюс почти 10k тестов.

Дальше --- авторегрессионный action space по схеме SAINT, RMCTS, пайплайн BC-PPO-Self-Play и opponent modeling.

\section*{Слайд 13 --- Выводы (20 сек)}

Резюмируя.

Первое --- headless-среда HS:BG с наиболее полным охватом механик по сравнению с Leiden и LUT.

Второе --- Transformer-экстрактор под автобаттлеры, новая сборка компонент из state-of-the-art работ 2018--2025 годов.

Третье --- экспериментально подтверждено ускорение в 5 раз по sample efficiency, плюс attention-анализ показал самоорганизованную специализацию голов.

\textbf{Спасибо за внимание.}

\hrulefill

\section*{Если надо сжать до 6 минут}

\begin{itemize}[noitemsep]
  \item Слайд 4 $\to$ 30 сек (выкини Fireplace/Unity, оставь только Leiden/LUT)
  \item Слайд 8 (PvP) $\to$ 20 сек
  \item Слайд 11 (ablation) $\to$ 25 сек
  \item \textbf{Не трогай} слайды 6 и 7 --- это ядро защиты
\end{itemize}

\section*{Критичные моменты}

\begin{itemize}[noitemsep]
  \item Говори <<\textbf{я} сделал>>, <<\textbf{моя} среда>>, <<\textbf{я} собрал>> --- чётко разграничивай своё и чужое (требование гайда Соколова)
  \item Для каждой чужой идеи называй фамилию + год: Perez 2018, Parisotto 2020, AlphaStar 2019, Set Transformer Lee 2019, MARL-GPT 2024
  \item На слайде 7 главная фраза --- \textbf{<<в 5 раз быстрее по sample efficiency>>} --- произнеси её громко и с паузой
  \item На слайде 11 не извиняйся за отсутствие полного ablation. Сформулируй как \textbf{ограничение вычислительных ресурсов} и \textbf{research trade-off}, обоснованный литобзором
  \item На слайдах 9--10 обязательно физически тыкай пальцем/курсором --- без этого attention-картинки мёртвые
\end{itemize}

\end{document}
